\subsection{4. The Core Workflow: Step-by-Step
Instructions}\label{the-core-workflow-step-by-step-instructions}

This is the exact sequence you must follow every time you work on a
replication.

\subsubsection{Phase 0: Starting a New
Session}\label{phase-0-starting-a-new-session}

Before doing anything else at the start of each work session:

\begin{minted}{bash}
# 1. Navigate to the repository
cd /Users/joaopms/Documents/Projeto_Mestrado

# 2. Check what's currently happening
git status
git log --oneline -5
\end{minted}

This tells you exactly where you left off.

\subsubsection{Phase 1: Picking an
Experiment}\label{phase-1-picking-an-experiment}

Open \texttt{README.md} in your browser or editor. The
Experiment Index table shows all 17 experiments and their current
status. Work in \textbf{priority order}:

\begin{itemize}
\tightlist
\item
  \textbf{P0 first} (001--005): Core thesis papers --- resampling
  strategies for imbalanced regression
\item
  \textbf{P1 second} (006--011): Methodology papers --- performance
  estimation, model selection
\item
  \textbf{P2 after} (012--015): Extended scope
\item
  \textbf{P3 last} (016--017): Reference implementations
\end{itemize}

Pick the experiment with the \textbf{lowest number} that is not yet
\texttt{[DONE] Done}.

\subsubsection{Phase 2: Setting Up the Experiment
Environment}\label{phase-2-setting-up-the-experiment-environment}

This is the most important and most manual phase. Do it in order.

\paragraph{Step 2.1 --- Read the paper and the
ARA}\label{step-2.1-read-the-paper-and-the-ara}

Before touching any code, read: 1. The experiment's
\texttt{README.md} --- understand the target results
you need to reproduce 2. The corresponding ARA in your Obsidian Vault
(if it exists) --- it contains distilled information about the paper's
methodology and code

\paragraph{\texorpdfstring{Step 2.2 --- Copy the original code into
\texttt{src/original/}}{Step 2.2 --- Copy the original code into src/original/}}\label{step-2.2-copy-the-original-code-into-srcoriginal}

This is \textbf{always done manually}. You have two options:

\textbf{Option A: Clone the repository}

\begin{minted}{bash}
cd experiments/NNN_folder_name/src/
git clone https://github.com/author/repo original
\end{minted}

\begin{quote}
\small {[}WARN{]} Do NOT use \texttt{git clone} inside the
main repo without careful \texttt{.gitignore} handling.
It's safer to clone elsewhere, then copy the folder.
\end{quote}

\textbf{Option B: Copy downloaded files}

\begin{minted}{bash}
# After downloading the zip and extracting:
cp -r ~/Downloads/author-repo/ experiments/NNN_folder_name/src/original/
\end{minted}

\paragraph{Step 2.3 --- Create an isolated Python
environment}\label{step-2.3-create-an-isolated-python-environment}

\textbf{CRITICAL: Each experiment must have its own isolated
environment.} Never install packages globally.

\begin{minted}{bash}
cd experiments/NNN_folder_name/

# Create the virtual environment
python3 -m venv .venv

# Activate it
source .venv/bin/activate

# Confirm you're in the right environment
which python3
# Should show: .../experiments/NNN_folder_name/.venv/bin/python3
\end{minted}

\paragraph{Step 2.4 --- Install
dependencies}\label{step-2.4-install-dependencies}

\begin{minted}{bash}
# Option A: If the paper has a requirements.txt
pip install -r src/original/requirements.txt

# Option B: If the paper has a setup.py or pyproject.toml
pip install -e src/original/

# Option C: Install manually (document every package you install)
pip install numpy pandas scikit-learn

# Save what you installed
pip freeze > requirements.txt
\end{minted}

\paragraph{Step 2.5 --- Test that everything imports without
error}\label{step-2.5-test-that-everything-imports-without-error}

\begin{minted}{bash}
# Quick smoke test — confirm no broken imports
python3 -c "import sys; sys.path.insert(0, 'src/original'); print('imports OK')"
\end{minted}

\paragraph{Step 2.6 --- Change status to {[}RUNNING{]} Setup and
commit}\label{step-2.6-change-status-to-running-setup-and-commit}

Once the environment works: 1. Open
\texttt{experiments/NNN\textbackslash\{\}\_folder\textbackslash\{\}\_name/README.md} 2.
Change
\texttt{\textbackslash\{\}\#\textbackslash\{\}\# Replication Status: [WAITING] Pending} to
\texttt{\textbackslash\{\}\#\textbackslash\{\}\# Replication Status: [RUNNING] Setup} 3.
Run the dashboard update:
\texttt{bash    python3 scripts/update\textbackslash\{\}\_readme.py} 4.
Commit:
\texttt{bash    git add experiments/NNN\textbackslash\{\}\_folder\textbackslash\{\}\_name/ README.md    git commit -m "[EXP-NNN] setup: clone original code and create venv"    git push origin main}

\subsubsection{Phase 3: Running the
Replication}\label{phase-3-running-the-replication}

\paragraph{Step 3.1 --- Identify the exact target
results}\label{step-3.1-identify-the-exact-target-results}

Read the paper carefully. Decide which specific tables and figures you
will reproduce. Add them to the experiment's
\texttt{README.md} under the
\texttt{Objective} section:

\begin{minted}{text}
## Objective

**Target results to replicate**:
- [ ] Table 3: RMSE comparison across resampling strategies on 5 datasets
- [ ] Figure 2: Performance vs. imbalance ratio curve
\end{minted}

\paragraph{Step 3.2 --- Run the original code with the original
parameters}\label{step-3.2-run-the-original-code-with-the-original-parameters}

Always start with the \textbf{exact original parameters} the authors
used. Do not change anything on the first run.

\begin{minted}{bash}
# Activate the environment first
source experiments/NNN_folder_name/.venv/bin/activate

# Navigate to the experiment directory
cd experiments/NNN_folder_name/

# Run the original code (the exact command varies by paper)
python3 src/original/main.py
# or: Rscript src/original/run_experiments.R
\end{minted}

\paragraph{Step 3.3 --- Update the REPLICATION\_LOG
immediately}\label{step-3.3-update-the-replication_log-immediately}

Every time you do something, write it down in
\texttt{REPLICATION\textbackslash\{\}\_LOG.md}. This is your scientific
diary. Entries go at the \textbf{top} of the file (most recent first).

\begin{minted}{text}
## [2026-07-15] First successful run attempt

**Time spent**: 2h 30m

### What was done
- Activated .venv and ran src/original/main.py
- Got a FileNotFoundError for the AEMO dataset

### What worked
- All imports resolved correctly
- Configuration loading worked

### Issues encountered
- The paper's data is not publicly available. Authors reference "private AEMO dataset".
- Contacted authors via email.

### Next steps
- Try with the public NASDAQ dataset as a proxy
- Wait for author response
\end{minted}

\paragraph{Step 3.4 --- Change status to {[}REVIEW{]} In
Progress}\label{step-3.4-change-status-to-review-in-progress}

Once you have your first successful run, update the status badge in the
experiment's \texttt{README.md} from
\texttt{[RUNNING]} to
\texttt{[REVIEW]}.

\subsubsection{Phase 4: Recording
Results}\label{phase-4-recording-results}

\paragraph{Step 4.1 --- Save results to the right
location}\label{step-4.1-save-results-to-the-right-location}

All outputs \textbf{must} go inside the experiment's
\texttt{results/} directory: - Tables -\textgreater{}
\texttt{results/tables/} as
\texttt{.csv} files - Figures -\textgreater{}
\texttt{results/figures/} as
\texttt{.png} (and optionally
\texttt{.pdf}) files

Name files descriptively:

\begin{minted}{text}
results/
├── tables/
│   ├── table3_rmse_comparison.csv
│   └── table4_mae_by_dataset.csv
└── figures/
    ├── figure2_performance_vs_imbalance.png
    └── figure5_training_time.png
\end{minted}

\paragraph{Step 4.2 --- Fill the comparison table in
README.md}\label{step-4.2-fill-the-comparison-table-in-readme.md}

This is the most important scientific output. For each metric you
reproduce:

\begin{minted}{text}
## Key Results

| Metric | Paper Reports | Our Replication | Δ (absolute) | Notes |
|--------|--------------|-----------------|--------------|-------|
| RMSE (AEMO dataset) | 0.342 | 0.349 | +0.007 | Different random seed |
| SERA (AEMO dataset) | 0.891 | 0.887 | -0.004 | Within acceptable range |
\end{minted}

\paragraph{Step 4.3 --- Document deviations
honestly}\label{step-4.3-document-deviations-honestly}

In the \texttt{Observations \textbackslash\{\}\& Deviations} section:

\begin{minted}{text}
## Observations & Deviations

- **Library version**: Original used scikit-learn 0.24.2; we used 1.3.0. 
  The API for `RandomForestRegressor` changed in v1.0; adjusted accordingly.
- **Dataset**: Original code references a private AEMO energy dataset. 
  We used the publicly available GEFCom2014 as a proxy. Results are directionally consistent.
- **Seed**: Original paper does not specify a seed. We used seed=42 for reproducibility.
\end{minted}

\subsubsection{Phase 5: Marking Complete and
Committing}\label{phase-5-marking-complete-and-committing}

\paragraph{Step 5.1 --- Mark checkboxes done in
README}\label{step-5.1-mark-checkboxes-done-in-readme}

Check off every target result you successfully reproduced:

\begin{minted}{text}
**Target results to replicate**:
- [x] Table 3: RMSE comparison across resampling strategies on 5 datasets
- [x] Figure 2: Performance vs. imbalance ratio curve
\end{minted}

\paragraph{Step 5.2 --- Change status to {[}DONE{]}
Done}\label{step-5.2-change-status-to-done-done}

Change
\texttt{\textbackslash\{\}\#\textbackslash\{\}\# Replication Status: [REVIEW] In Progress}
to \texttt{\textbackslash\{\}\#\textbackslash\{\}\# Replication Status: [DONE] Done}.

\paragraph{Step 5.3 --- Run the update script and
commit}\label{step-5.3-run-the-update-script-and-commit}

\begin{minted}{bash}
# Update the main dashboard
python3 scripts/update_readme.py

# Stage all changes
git add experiments/NNN_folder_name/ README.md

# Commit with a descriptive message
git commit -m "[EXP-NNN] results: reproduce Table 3 and Figure 2"

# Push
git push origin main
\end{minted}
